\scanpage{77}
II Cohomologie de Weil \unsure{\ldots} [\unsure{semble marcher moins bien}\ldots]

Il n'est plus vrai en général que \inserted{\unsure{Supposons} $A$ \uline{\unsure{hensél.}}}
\[
H^i(A,G)=0\qquad\text{si }i>0.
\]
\deleted{\unsure{Pour}} Si $G$ est $k$-\uline{\unsure{simple|ample}}, alors $H^i(A,G)=H^i(L,G)$ \inserted{si $i>0$}, $L$ étant le corps résiduel, et \deleted{si $i>0$} on a
\[
H^i(L,G)=H^i(L,G/\unsure{\mathscr L|\alpha})
\]
où $\mathscr L$ est le plus grand ss-gpe linéaire connexe de $G$.
\deleted{$\unsure{H_A}$ \unsure{ne serait pas} \illegible{2}}
\deleted{$\unsure{H_A}$ \unsure{ce qui est} \illegible{2}}

\deleted{$H^1(X/X-a,G)\unsure{-|\longrightarrow}H^1(X,G)$ \unsure{à déterminer à travers les}}
\deleted{Si $\dim A=0$, rien de nouveau.}
\deleted{Si $\dim A=1$, on a $H^1_{\unsure{l.t.|l.tr.}}(X,G)$, \unsure{donc est nul, d'où on aura}}
\[
\deleted{H^1_{\unsure{K|A}}(G)=\operatorname{Coker}(H^0(X)\unsure{-|\longrightarrow}H^0(X-a))}
\]

Si $\dim A=0$, rien de nouveau.

Si $\dim A\geq1$, alors $H^1(X/X-a,G)$ est extension de $\operatorname{Ker}(H^1(X,\unsure{G})-H^1(X-a,G))$ par $G(K)/G(A)$.

\inserted{Or} \deleted{\unsure{on a}}
\[
\operatorname{Ker}(H^1(X,G)\longrightarrow H^1(X-a,G))=0
\]
% lower-left marginal diagram
\margin{\diagram{\begin{tikzcd}[ampersand replacement=\&,column sep=small,row sep=small]
0\arrow[d]\&\\
H^1(X,G)\arrow[r]\arrow[d]\&H^1(X-a,G)\arrow[d]\\
H^1(X,G/\mathscr L)\arrow[r]\&H^1(X-a,G/\mathscr L)
\end{tikzcd}}}
si $G$ est \unsure{simple|ample}, car \deleted{\unsure{une section de}} \inserted{\unsure{on est ramené au}} cas où $G$ est extension d'un grp discret par un grp propre, mais alors une section \deleted{\unsure{d'un}} rationnelle d'un fibré principal est partout définie [car définie en codim $1$, donc partout].
% archive sheet number: 76
